% Part: first-order-logic
% Chapter: natural-deduction
% Section: proving-things

\documentclass[../../../include/open-logic-section]{subfiles}

\begin{document}

\iftag{FOL}
      {\olfileid[hi]{fol}{ntd}{pro}}
      {\olfileid[hi]{pl}{ntd}{pro}}

\olsection{\usetoken{p}{derivation} बनाने के उदाहरण}

\begin{ex}
आइए !!a{derivation} दें जिसका निष्कर्ष !!{sentence} $(!A \land !B) \lif !A$
है।

वांछित निष्कर्ष को !!{derivation} के सबसे नीचे लिखकर आरम्भ करते हैं।
\begin{prooftree}
\AxiomC{}
\UnaryInfC{$(!A\land !B) \lif !A$}
\end{prooftree}

अब तय करना है कि किस प्रकार का अनुमान इस रूप का !!a{sentence} दे सकता है।
निष्कर्ष का !!{main operator} $\lif$ है, इसलिए \Intro{\lif} नियम से निष्कर्ष
तक पहुँचने का प्रयत्न करेंगे। आगे बढ़ते हुए संबद्ध मान्यताओं को लिखना और
अनुमान-नियमों पर लेबल लगाना उचित है; इससे प्रमाण के अंत में आसानी से देख
सकते हैं कि सभी मान्यताएँ !!{discharged} हुई हैं या नहीं।
\begin{prooftree}
\AxiomC{$\Discharge{!A \land !B}{1}$}
\DeduceC{$!A$}
\DischargeRule{\Intro{\lif}}{1} 
\UnaryInfC{$(!A\land !B) \lif !A$}
\end{prooftree}

अब मान्यता $!A \land !B$ से $!A$ तक के चरण भरने हैं। यहाँ केवल एक संयोजक
$\land$ है, इसलिए $\land$ निरसन-नियम लगाना होगा। इससे यह प्रमाण मिलता है:
\begin{prooftree}
\AxiomC{$\Discharge{!A \land !B}{1}$}
\RightLabel{\Elim{\land}}
\UnaryInfC{$!A$}
\DischargeRule{\Intro{\lif}}{1} 
\UnaryInfC{$(!A\land !B) \lif !A$}
\end{prooftree}
अब हमारे पास $(!A \land
!B) \lif !A$ की सही !!{derivation} है।
\end{ex}

\begin{ex}
अब !!a{derivation} दें जिसका निष्कर्ष $(\lnot !A \lor !B)
\lif (!A \lif !B)$ है।

वांछित निष्कर्ष को !!{derivation} के सबसे नीचे लिखकर आरम्भ करते हैं।
\begin{prooftree}
\AxiomC{}
\UnaryInfC{$(\lnot !A \lor !B) \lif (!A \lif !B)$}
\end{prooftree}
यह निष्कर्ष देने वाला तार्किक नियम खोजने के लिए निष्कर्ष के तार्किक संयोजक
देखते हैं: $\lnot$, $\lor$, और $\lif$। अभी केवल $\lif$ की पहली आवृत्ति पर
ध्यान देंगे, क्योंकि वही !!{main operator} है—अंतिम अनुक्रम के !!{sentence}
का; $\lnot$, $\lor$, और $\lif$ की दूसरी आवृत्ति किसी अन्य संयोजक के परास में
हैं, इसलिए उन्हें बाद में सँभालेंगे। अतः \Intro{\lif} नियम से आरम्भ करते हैं।
उसका सही प्रयोग इस रूप का होगा:
\begin{prooftree}
\AxiomC{$\Discharge{\lnot !A \lor !B}{1}$}
\DeduceC{$!A \lif !B$}
\DischargeRule{\Intro{\lif}}{1}
\UnaryInfC{$(\lnot !A \lor !B) \lif (!A \lif !B)$}
\end{prooftree}
अब आगे बढ़ने के दो मार्ग हैं। या तो नीचे से ऊपर काम जारी रखकर
\Intro{\lif} का एक और प्रयोग खोजें, या ऊपर से नीचे काम करके \Elim{\lor} नियम
लगाएँ। दूसरा मार्ग लेते हैं। \Elim{\lor} के सबसे बाएँ पूर्वाधार के रूप में
मान्यता $\lnot !A \lor !B$ लेंगे। \Elim{\lor} के सही प्रयोग में अन्य दोनों
पूर्वाधार निष्कर्ष $!A \lif !B$ के समान होने चाहिए, पर उनमें से प्रत्येक को
बारी-बारी से एक अन्य मान्यता—अर्थात् $\lnot !A \lor !B$ के दो वियोज्यों में
से एक—से व्युत्पन्न किया जा सकता है। इसलिए हमारी !!{derivation} ऐसी होगी:
\begin{prooftree}
\AxiomC{$\Discharge{\lnot !A \lor !B}{1}$}
\AxiomC{$\Discharge{\lnot !A}{2}$}
\DeduceC{$!A \lif !B$}
\AxiomC{$\Discharge{!B}{2}$}
\DeduceC{$!A \lif !B$}
\DischargeRule{\Elim{\lor}}{2}
\TrinaryInfC{$!A \lif !B$}
\DischargeRule{\Intro{\lif}}{1} 
\UnaryInfC{$(\lnot !A \lor !B) \lif (!A \lif !B)$}
\end{prooftree}

दाईं ओर की दोनों शाखाओं में हम !!{derive} करना चाहते हैं $!A
\lif !B$ को; इसके लिए \Intro{\lif} लगाना सबसे उचित है।
\begin{prooftree}
\AxiomC{$\Discharge{\lnot !A \lor !B}{1}$}
\AxiomC{$\Discharge{\lnot !A}{2}, \Discharge{!A}{3}$}
\DeduceC{$!B$}
\DischargeRule{\Intro{\lif}}{3}
\UnaryInfC{$!A \lif !B$}
\AxiomC{$\Discharge{!B}{2}, \Discharge{!A}{4}$}
\DeduceC{$!B$}
\DischargeRule{\Intro{\lif}}{4}
\UnaryInfC{$!A \lif !B$}
\DischargeRule{\Elim{\lor}}{2}
\TrinaryInfC{$!A \lif !B$}
\DischargeRule{\Intro{\lif}}{1} 
\UnaryInfC{$(\lnot !A \lor !B) \lif (!A \lif !B)$}
\end{prooftree}

पूरी !!{derivation} के दो अधूरे भागों में हमें !!{derivation}s चाहिए: $!B$ की—
बीच में $\lnot !A$ और $!A$ से तथा बाईं ओर $!A$ और $!B$ से। पहले बीच वाली
लेते हैं। $\lnot
!A$ और $!A$, \Elim{\lnot} के दो पूर्वाधार हैं:
\begin{prooftree}
\AxiomC{$\Discharge{\lnot !A}{2}$}
\AxiomC{$\Discharge{!A}{3}$}
\RightLabel{\Elim{\lnot}}
\BinaryInfC{$\lfalse$}
\DeduceC{$!B$}
\end{prooftree}
\FalseInt का प्रयोग करके $!B$ को निष्कर्ष के रूप में पा सकते हैं और शाखा
पूरी कर सकते हैं।
\begin{prooftree}
\AxiomC{$\Discharge{\lnot !A \lor !B}{1}$}
\AxiomC{$\Discharge{\lnot !A}{2}$}
\AxiomC{$\Discharge{!A}{3}$}
\RightLabel{\Intro{\lfalse}}
\BinaryInfC{$\lfalse$}
\RightLabel{\FalseInt}
\UnaryInfC{$!B$}
\DischargeRule{\Intro{\lif}}{3}
\UnaryInfC{$!A \lif !B$}
\AxiomC{$\Discharge{!B}{2}, \Discharge{!A}{4}$}
\DeduceC{$!B$}
\DischargeRule{\Intro{\lif}}{4}
\UnaryInfC{$!A \lif !B$}
\DischargeRule{\Elim{\lor}}{2}
\TrinaryInfC{$!A \lif !B$}
\DischargeRule{\Intro{\lif}}{1} 
\UnaryInfC{$(\lnot !A \lor !B) \lif (!A \lif !B)$}
\end{prooftree}

अब सबसे दाहिनी शाखा देखें। यहाँ यह समझना महत्त्वपूर्ण है कि !!{derivation}
की परिभाषा मान्यताओं को \emph{मुक्त करने की अनुमति देती है}, पर ऐसा करना
\emph{अनिवार्य नहीं} बनाती। अर्थात् यदि $!B$ को मान्यताओं $!A$ और $!B$ में से
केवल एक का उपयोग करके व्युत्पन्न कर सकें, तो यह स्वीकार्य है। और !!{derive}
करना $!B$ को~$!B$ से सर्वथा सरल है: $!B$ स्वयं ऐसी !!a{derivation} है और
किसी अनुमान की आवश्यकता नहीं। अतः मान्यता~$!A$ को सीधे हटा सकते हैं।
\begin{prooftree}
\AxiomC{$\Discharge{\lnot !A \lor !B}{1}$}
\AxiomC{$\Discharge{\lnot !A}{2}$}
\AxiomC{$\Discharge{!A}{3}$}
\RightLabel{\Elim{\lnot}}
\BinaryInfC{$\lfalse$}
\RightLabel{\FalseInt}
\UnaryInfC{$!B$}
\DischargeRule{\Intro{\lif}}{3}
\UnaryInfC{$!A \lif !B$}
\AxiomC{$\Discharge{!B}{2}$}
\RightLabel{\Intro{\lif}}
\UnaryInfC{$!A \lif !B$}
\DischargeRule{\Elim{\lor}}{2}
\TrinaryInfC{$!A \lif !B$}
\DischargeRule{\Intro{\lif}}{1} 
\UnaryInfC{$(\lnot !A \lor !B) \lif (!A \lif !B)$}
\end{prooftree}
ध्यान दें कि पूर्ण !!{derivation} में सबसे दाहिना \Intro{\lif} अनुमान वास्तव
में किसी मान्यता को मुक्त नहीं करता।
\end{ex}

\begin{ex}
अब तक हमें \FalseCl{} नियम की आवश्यकता नहीं पड़ी। इसकी विशेषता यह है कि यह
ऐसी मान्यता को मुक्त करने देता है जो नियम के निष्कर्ष का उप-!!{formula} नहीं
है। इसका \FalseInt{} नियम से निकट संबंध है। वास्तव में \FalseInt{} नियम,
\FalseCl{} नियम की एक विशेष दशा है---``अंतःप्रज्ञावादी तर्कशास्त्र'' नाम का
एक तर्कशास्त्र है जिसमें केवल \FalseInt{} अनुमत है। जब कोई और उपाय काम न करे,
तब \FalseCl{} अंतिम उपाय है। उदाहरणतः, मान लें कि हमें !!{derive} करना है
$!A \lor \lnot !A$। सामान्य रणनीति $!A \lor \lnot !A$ को $\Intro{\lor}$
लगाकर !!{derive} करने की होगी। पर इसके लिए बिना किसी मान्यता के !!{derive} करना
होगा या तो $!A$ को या $\lnot !A$ को, और यह सम्भव नहीं। यहाँ \FalseCl{} काम
आता है!
\begin{prooftree}
  \AxiomC{$\Discharge{\lnot(!A \lor \lnot !A)}{1}$}
  \DeduceC{$\lfalse$}
  \DischargeRule{\FalseCl}{1}
  \UnaryInfC{$!A \lor \lnot !A$}
\end{prooftree}
अब !!a{derivation} खोज रहे हैं: $\lfalse$ की, $\lnot(!A \lor
\lnot !A)$ से। चूँकि $\lfalse$, $\Elim{\lnot}$ का निष्कर्ष है, इसे लगाने का
प्रयत्न कर सकते हैं:
\begin{prooftree}
  \AxiomC{$\Discharge{\lnot(!A \lor \lnot !A)}{1}$}
  \DeduceC{$\lnot !A$}
  \AxiomC{$\Discharge{\lnot(!A \lor \lnot !A)}{1}$}
  \DeduceC{$!A$}
  \RightLabel{\Elim{\lnot}}
  \BinaryInfC{$\lfalse$}
  \DischargeRule{\FalseCl}{1}
  \UnaryInfC{$!A \lor \lnot !A$}
\end{prooftree}
हमारी रणनीति में !!a{derivation} खोजनी है:~$\lnot !A$ की; इसके लिए
~$\Intro{\lnot}$ का प्रयोग चाहिए:
\begin{prooftree}
  \AxiomC{$\Discharge{\lnot(!A \lor \lnot !A)}{1}, \Discharge{!A}{2}$}
  \DeduceC{$\lfalse$}
  \DischargeRule{\Intro{\lnot}}{2}
  \UnaryInfC{$\lnot !A$}
  \AxiomC{$\Discharge{\lnot(!A \lor \lnot !A)}{1}$}
  \DeduceC{$!A$}
  \RightLabel{\Elim{\lnot}}
  \BinaryInfC{$\lfalse$}
  \DischargeRule{\FalseCl}{1}
  \UnaryInfC{$!A \lor \lnot !A$}
\end{prooftree}
यहाँ $\lfalse$ को $\Elim{\lnot}$ लगाकर सरलता से पा सकते हैं—मान्यता
$\lnot(!A \lor \lnot !A)$ और $!A \lor \lnot !A$ पर; दूसरा वाक्य हमारी नई
मान्यता $!A$ से~$\Intro{\lor}$ द्वारा मिलता है:
\begin{prooftree}
  \AxiomC{$\Discharge{\lnot(!A \lor \lnot !A)}{1}$}
  \AxiomC{$\Discharge{!A}{2}$}
  \RightLabel{\Intro{\lor}}
  \UnaryInfC{$!A \lor \lnot !A$}
  \RightLabel{\Elim{\lnot}}
  \BinaryInfC{$\lfalse$}
  \DischargeRule{\Intro{\lnot}}{2}
  \UnaryInfC{$\lnot !A$}
  \AxiomC{$\Discharge{\lnot(!A \lor \lnot !A)}{1}$}
  \DeduceC{$!A$}
  \RightLabel{\Elim{\lnot}}
  \BinaryInfC{$\lfalse$}
  \DischargeRule{\FalseCl}{1}
  \UnaryInfC{$!A \lor \lnot !A$}
\end{prooftree}
दाईं ओर वही रणनीति अपनाते हैं, केवल $!A$ को~\FalseCl से पाते हैं:
\begin{prooftree}
  \AxiomC{$\Discharge{\lnot(!A \lor \lnot !A)}{1}$}
  \AxiomC{$\Discharge{!A}{2}$}
  \RightLabel{\Intro{\lor}}
  \UnaryInfC{$!A \lor \lnot !A$}
  \RightLabel{\Elim{\lnot}}
  \BinaryInfC{$\lfalse$}
  \DischargeRule{\Intro{\lnot}}{2}
  \UnaryInfC{$\lnot !A$}
  \AxiomC{$\Discharge{\lnot(!A \lor \lnot !A)}{1}$}
  \AxiomC{$\Discharge{\lnot !A}{3}$}
  \RightLabel{\Intro{\lor}}
  \UnaryInfC{$!A \lor \lnot !A$}
  \RightLabel{\Elim{\lnot}}
  \BinaryInfC{$\lfalse$}
  \DischargeRule{\FalseCl}{3}
  \UnaryInfC{$!A$}
  \RightLabel{\Elim{\lnot}}
  \BinaryInfC{$\lfalse$}
  \DischargeRule{\FalseCl}{1}
  \UnaryInfC{$!A \lor \lnot !A$}
\end{prooftree}
\end{ex}

\begin{prob}
निम्नलिखित को दिखाने वाली !!{derivation}s दीजिए:
\begin{enumerate}
\item $!A \land (!B \land !C) \Proves (!A \land !B) \land !C$.
\item $!A \lor (!B \lor !C) \Proves (!A \lor !B) \lor !C$.
\item $!A \lif (!B \lif !C) \Proves !B \lif (!A \lif !C)$.
\item $!A \Proves \lnot\lnot !A$.
\end{enumerate}
\end{prob}

\begin{prob}
निम्नलिखित को दिखाने वाली !!{derivation}s दीजिए:
\begin{enumerate}
\item $(!A \lor !B) \lif !C \Proves !A \lif !C$.
\item $(!A \lif !C) \land (!B \lif !C) \Proves (!A \lor !B) \lif !C$.
\item $\Proves \lnot(!A \land \lnot !A)$.
\item $!B \lif !A \Proves \lnot !A \lif \lnot !B$.
\item $\Proves (!A \lif \lnot !A) \lif \lnot !A$.
\item $\Proves \lnot(!A \lif !B) \lif \lnot !B$.
\item $!A \lif !C \Proves \lnot (!A \land \lnot !C)$.
\item $!A \land \lnot !C \Proves \lnot (!A \lif !C)$.
\item $!A \lor !B, \lnot !B \Proves !A$.
\item $\lnot !A \lor \lnot !B \Proves \lnot(!A \land !B)$.
\item $\Proves (\lnot !A \land \lnot !B) \lif\lnot(!A \lor !B)$.
\item $\Proves \lnot(!A \lor !B) \lif (\lnot !A \land \lnot !B)$.
\end{enumerate}
\end{prob}

\begin{prob}
निम्नलिखित को दिखाने वाली !!{derivation}s दीजिए:
\begin{enumerate}
\item $\lnot(!A \lif !B) \Proves !A$.
\item $\lnot(!A \land !B) \Proves \lnot !A \lor \lnot !B$.
\item $!A \lif !B \Proves \lnot !A \lor !B$.
\item $\Proves \lnot \lnot !A \lif !A$.
\item $!A \lif !B, \lnot !A \lif !B \Proves !B$.
\item $(!A \land !B) \lif !C \Proves (!A \lif !C) \lor (!B \lif !C)$.
\item $(!A \lif !B) \lif !A \Proves !A$.
\item $\Proves (!A \lif !B) \lor (!B \lif !C)$.
\end{enumerate}
(इन सभी के लिए $\FalseCl$~नियम आवश्यक है।)
\end{prob}

\end{document}
